\chapter{保护神经系统}\label{ch:g8:protecting-the-nervous-system}

前两章已经说明\omterm{def:g5:brain-in-command:system}{神经系统}是什么：一套布线，其连接
会学习、会做决定，也会被欺骗
（\cref{rem:g8:nervous-communication:vulnerable}）。
本学年的结尾，把这些结构知识转化为保护。
\omterm{prop:g1:healthy-habits:needs}{睡眠}不足、噪声、屏幕以及精神活性物质，
究竟会在每一个连接上对你这台机器做什么？
在实际生活中，保护它是什么样子，
对十四岁的你又意味着什么？

\section{维护：睡眠与感官的极限}

\begin{proposition}[睡眠即突触维护]\label{prop:g8:protecting-the-nervous-system:sleep}
\cref{prop:g5:brain-in-command:protect}所说的归档与修复，
可以在连接层面来理解：\omterm{prop:g1:healthy-habits:needs}{睡眠}中，白天得到增强的\omterm{def:g8:nervous-communication:synapse}{突触}
被巩固——熟能生巧
（\cref{ex:g8:nervous-communication:juggler}）很大程度上
发生在夜间；\omterm{def:g5:brain-in-command:system}{大脑}的日常清理
也会清除白天积累的化学损耗。\omterm{prop:g1:healthy-habits:needs}{睡眠}不足时，昨天的学习
只被半归档，明天的判断则变得迟缓：青少年的\omterm{def:g5:brain-in-command:system}{大脑}
正在重新布线（\cref{rem:g8:puberty:moods}），
\omterm{def:g6:exploring-our-environment:components}{生物}钟又偏向晚睡，
恰恰在学校开始得最早的时候，
最需要\omterm{prop:g1:healthy-habits:needs}{睡眠}的时间。
\end{proposition}
\admitted

\begin{example}[复习的悖论]\label{ex:g8:protecting-the-nervous-system:revision}
考试前，两名学生各自复习：一人复习到午夜，又继续到凌晨两点；
另一人复习到十点便去睡觉。
通常是后者成绩更好——材料已经穿过他们的连接，
并且\emph{得到了巩固}；前者多出来的时间，
则损失了让复习真正留存下来的那一步归档。
“睡一觉再说”，其实就是\omterm{def:g8:nervous-communication:synapse}{突触}层面的工程。
\end{example}

\begin{proposition}[声音过响即是剂量]\label{prop:g8:protecting-the-nervous-system:noise}
耳朵中的感觉\omterm{def:g5:first-look-at-cells:cell}{细胞}——那些将声音转化为\omterm{def:g8:nervous-communication:neuron}{神经信息}的细微感受器——
不会再生（\cref{prop:g1:the-five-senses:fragile}曾温和地提到这一点，
这里给出其机制）。强声会按剂量杀伤它们：
音量乘以暴露的时间。音乐会和以大音量使用的耳机，
逐小时消耗着一笔有限的感受器储备，
而\omterm{prop:g1:healthy-habits:needs}{睡眠}无法将其\omterm{def:g7:organs-at-work:recovery}{恢复}——一夜喧闹之后耳中的嗡鸣
（\cref{ex:g1:the-five-senses:loud}）正是过量的警示。
规则是一种预算：音量调低，
时间缩短，中间多休息；
在扩音设备占据主导的场合，则要戴上耳塞。
\end{proposition}
\admitted

\section{令连接失常的物质}

\begin{proposition}[精神活性物质作用于突触]\label{prop:g8:protecting-the-nervous-system:substances}
凡是改变情绪、感知或警觉度的物质，
都作用于\cref{def:g8:nervous-communication:synapse}中的间隙——
或模拟某种信使，或阻断某种信使，
或迫使其释放：
\begin{itemize}
  \item \emph{酒精}会对整个系统的连接起阻尼作用：
        回路从头到尾都变慢——判断力、\omterm{def:g5:brain-in-command:reaction}{反应时间}
        （\cref{ex:g5:brain-in-command:driver}中以米计的刹车距离会变长），
        平衡与记忆也会受影响；仍在发育的青少年\omterm{def:g5:brain-in-command:system}{大脑}，
        其尚未建成的连接，
        受到的阻尼作用最深；
  \item \emph{烟草中的尼古丁}模拟一种天然信使，
        反复作用之后，会把奖赏连接改造成趋向
        \emph{\omterm{prop:g5:healthy-choices:substances}{成瘾}}的方向
        （\cref{prop:g5:healthy-choices:substances}所说的陷阱，
        此处给出其机制：\omterm{def:g8:nervous-communication:synapse}{突触}会调整以\omterm{def:g4:adapted-to-their-world:adaptation}{适应}这一剂量，
        没有它时便会“抗议”）；
  \item \emph{大麻}模拟另一类信使家族，
        搅乱的正是记忆、动机与时序方面的连接——
        同样地，代价最高的
        是仍在构建中的\omterm{def:g5:brain-in-command:system}{大脑}；
  \item 物质越烈，对连接的强行干扰就越粗暴，
        而奖赏重布线也越快，
        正是\omterm{prop:g5:healthy-choices:substances}{成瘾}的引擎。
\end{itemize}
\end{proposition}
\admitted

\begin{remark}[成瘾的机制]\label{rem:g8:protecting-the-nervous-system:addiction}
\Cref{rem:g5:healthy-choices:never}所说的陷阱，如今有了线路图：
奖赏连接在反复剂量作用下，
\emph{\omterm{def:g4:adapted-to-their-world:adaptation}{适应}}——受体减少，阈值升高——
直到需要该物质才能感觉正常，
一旦缺失，系统便会发出“抗议”。
这种\omterm{def:g4:adapted-to-their-world:adaptation}{适应}，正是“进入容易、退出困难”
这一不对称性在\omterm{def:g8:nervous-communication:synapse}{突触}层面的体现：开始是一个决定；
停止则要对抗已经重新布线的连接。
仍具可塑性的青少年\omterm{def:g5:brain-in-command:system}{大脑}重布线最快，
这也是为什么各项统计都显示，早期开始的人陷得最深，
而“永不开始”始终是最省事的选择。
\end{remark}

\begin{example}[屏幕的循环]\label{ex:g8:protecting-the-nervous-system:screens}
没有任何分子参与，但连接同样受到作用：信息流和游戏
经过刻意设计——用不可预测的小奖励撩拨奖赏线路，
这种安排最能训练连接；
结果是\omterm{prop:g5:healthy-choices:substances}{成瘾}陷阱的一个温和“远亲”，
表现为刷屏难以停下、夜晚悄然消失，
\omterm{prop:g1:healthy-habits:needs}{睡眠}被推迟（\cref{prop:g5:healthy-choices:screens}中的时长，
加上\cref{prop:g8:protecting-the-nervous-system:sleep}中的代价）。
防御的办法一如既往：事先定好使用的时长，
并使
最后一小时远离屏幕。
\end{example}

\section{实践：自己落实保护}

\begin{method}[神经系统使用说明书]\label{met:g8:protecting-the-nervous-system:manual}
\cref{met:g5:healthy-choices:run}的护理清单，
附上机制后重新列出：
\begin{enumerate}
  \item \emph{保证\omterm{prop:g1:healthy-habits:needs}{睡眠}时长}——巩固与清理在夜间进行；
        守住最后一小时，
        远离屏幕；
  \item \emph{控制响度预算}——感受器的
        本钱不会再生：调低音量，留出停歇，
        噪声轰鸣处戴上耳塞；
  \item \emph{戴头盔}，用于跌落迅速的场合
        （\cref{ex:g5:brain-in-command:helmet}）——线路
        不会像骨头那样愈合；
  \item \emph{不碰骗过连接的东西}，适用于
        尚在发育中的\omterm{def:g5:brain-in-command:system}{大脑}——“不”要冷静地
        提前定好；
        （\cref{ex:g5:healthy-choices:answers}中的随身那句话，
        现附完整理由）；
  \item \emph{善用\omterm{def:g5:brain-in-command:system}{大脑}这台机器}——
        学习、运动、音乐：连接随使用而增强，
        训练有素的回路，是其他习惯所
        依托的资本。
\end{enumerate}
\end{method}

\begin{example}[求助有门，且行之有效]\label{ex:g8:protecting-the-nervous-system:help}
对于已经陷进去的人——无论是自己的习惯还是朋友的，
无论涉及物质还是屏幕：\cref{ex:g8:contraception:help}中的
求助联系方式同样适用于此。学校护士和医生
按本章所示的那样看待\omterm{prop:g5:healthy-choices:substances}{成瘾}——一种\omterm{def:g8:nervous-communication:neuron}{神经}连接状况，
而非品行判定，
已重新连接的\omterm{def:g8:nervous-communication:synapse}{突触}可以重新\omterm{def:g4:adapted-to-their-world:adaptation}{适应}：
比从未沾染更难，但远非毫无希望。最糟的办法
是操场上的那一套：沉默。
\end{example}

\section{练习}

\begin{exercise}[$\star$]\label{exo:g8:protecting-the-nervous-system:1}
\omterm{prop:g1:healthy-habits:needs}{睡眠}在连接处起什么作用，
谁最需要这些\omterm{prop:g1:healthy-habits:needs}{睡眠}时长？
\end{exercise}

\begin{exercise}[$\star$]\label{exo:g8:protecting-the-nervous-system:2}
为什么大声是一个“预算”问题？花掉的是什么，以及
过量的警示是什么？
\end{exercise}

\begin{exercise}[$\star$]\label{exo:g8:protecting-the-nervous-system:3}
所有精神活性物质作用于何处，
以哪三种方式起作用？
\end{exercise}

\begin{exercise}[$\star$]\label{exo:g8:protecting-the-nervous-system:4}
酒精对这一回路有何影响——请列出三个
环节。
\end{exercise}

\begin{exercise}[$\star$]\label{exo:g8:protecting-the-nervous-system:5}
用两句话陈述\omterm{prop:g5:healthy-choices:substances}{成瘾}的机制。
\end{exercise}

\begin{exercise}[$\star$]\label{exo:g8:protecting-the-nervous-system:6}
屏幕回路为何属于本章，
尽管它不涉及分子？
\end{exercise}

\begin{exercise}[$\star\star$]\label{exo:g8:protecting-the-nervous-system:7}
借助\cref{prop:g8:protecting-the-nervous-system:sleep}，
解释“复习悖论”。
\end{exercise}

\begin{exercise}[$\star\star$]\label{exo:g8:protecting-the-nervous-system:8}
为什么相同剂量给青少年\omterm{def:g5:brain-in-command:system}{大脑}带来的代价比成人更大？
请给出两种机制
（\cref{rem:g8:puberty:moods}、
\cref{rem:g8:protecting-the-nervous-system:addiction}）。
\end{exercise}

\begin{exercise}[$\star\star$]\label{exo:g8:protecting-the-nervous-system:9}
将\cref{ex:g5:brain-in-command:driver}的米数
与酒精的阻尼相联系：延长的是什么，
为什么酒驾法规本质上是\omterm{def:g5:brain-in-command:reaction}{反应时间}的法规？
\end{exercise}

\begin{exercise}[$\star\star$]\label{exo:g8:protecting-the-nervous-system:10}
用连接层面的机制，重新论证使用说明书中的
三条内容。
\end{exercise}

\begin{exercise}[$\star\star$]\label{exo:g8:protecting-the-nervous-system:11}
“一种\omterm{def:g8:nervous-communication:neuron}{神经}连接状况，而非品行判定。”
这一重新界定在寻求帮助上有何改变，
以及在向朋友提供帮助时有何改变？
\end{exercise}

\begin{exercise}[$\star\star\star$]\label{exo:g8:protecting-the-nervous-system:12}
写出本学年的总结论证：依据
\cref{def:g8:nervous-communication:synapse}中可调的间隙，
推出\omterm{def:g5:brain-in-command:system}{大脑}最大的长处与其特有的脆弱之处——
学习与\omterm{prop:g5:healthy-choices:substances}{成瘾}，
即同一机制的两种读法。
\end{exercise}

\section{问题：预防之夜}

\begin{problem}[{周末问题——学校的家长与学生之夜，
依托机制设计的
脚本}]\label{pb:g8:protecting-the-nervous-system:1}
全班筹备学校的预防之夜：四个展位——
\omterm{prop:g1:healthy-habits:needs}{睡眠}、声音、物质、屏幕——每个展位
各配一项演示和一张机制卡片。

\medskip\noindent\textbf{第一部分——展位的
卡片。}
\begin{enumerate}
  \item 写出\omterm{prop:g1:healthy-habits:needs}{睡眠}卡片：演示用两名志愿者的
        标尺测试成绩，
        （\cref{met:g5:brain-in-command:ruler}），比较休息充足组与\omterm{prop:g1:healthy-habits:needs}{睡眠}不足组；
        预测两组成绩，并
        给出连接机制。
  \item 写出声音卡片：分贝计、一副
        耳塞，以及一句“你无法让它们
        再长回来”。补全机制。
  \item 写出物质卡片：一幅
        \omterm{def:g8:nervous-communication:synapse}{突触}示意图，三支箭头分别标注“模仿”“阻断”“促使释放”，
        将酒精、尼古丁与
        \cref{exo:g8:nervous-communication:11}中的一种毒液，
        分别配到对应箭头上。
  \item 写出屏幕卡片：没有分子参与，那么
        这个展位的带通知提醒的手机演示，
        与物质
        展位有何共同之处？
\end{enumerate}

\medskip\noindent\textbf{第二部分——来自现场的
提问。}
\begin{enumerate}[resume]
  \item 一位家长问：“我们这个年纪也喝过酒，
        不也好好的吗？”请结合工地上的不对称来回答：
        十四岁时有何不同？
  \item 一位学生说：“一支烟
        改变不了什么。”请结合
        \cref{rem:g8:protecting-the-nervous-system:addiction}
        以及\cref{rem:g5:healthy-choices:never}，第一剂的
        真实代价是什么？
  \item 一位家长问：“屏幕不就是新的
        电视吗？”请区分两种强化程序：
        什么使得不可预测的小奖励
        成为更难对付的“训练者”？
  \item 一位学生轻声问：“如果有人已经
        停不下来呢？”请给出
        \cref{ex:g8:protecting-the-nervous-system:help}的
        答案，包括求助地址。
\end{enumerate}

\medskip\noindent\textbf{第三部分——结束
讲话。}
\begin{enumerate}[resume]
  \item 讲话以杂耍者开头，以
        \omterm{prop:g5:healthy-choices:substances}{成瘾}结尾——“同一种可调节性，读两遍”。
        请用四句话写出这一论证。
  \item 随后，每位学生将拿到
        \cref{ex:g5:healthy-choices:answers}中的随身语句。解释
        包括机制在内，说明为什么这句话必须
        \emph{现在}就编写。
  \item 当晚的海报展示用户手册上的
        五条内容。请按十四岁青少年的
        一周安排排序，保护作用最强者居前，并
        为你的第一选择辩护。
  \item 用一句话作结：所保护的是什么，
        为什么十四岁是决定性十年的
        中段？
\end{enumerate}
\end{problem}
